%!TEX root = ../main.tex
% =============================================================
%  CAPÍTULO 2 — METODOLOGÍA             (Entrega 2 · 5 % · 5/5)
% =============================================================

\chapter{Metodología y procedimientos}\label{cap:metodologia}

\begin{guia}[Guía de la cátedra: qué se evalúa en este capítulo]
La cátedra exige \textbf{dos metodologías}. El error más común es
entregar solo una.
\begin{enumerate}
    \item \textbf{Metodología de la investigación}: área temática,
          planteamiento y delimitación del problema, marco teórico
          previsto, diseño concreto (fases), indicadores medibles con
          fuente y umbral, técnicas de recolección, instrumentos, datos,
          procesamiento, análisis crítico y síntesis. Tipo de
          investigación: en general, aplicada.
    \item \textbf{Metodología de dirección del proyecto}: marco de trabajo
          justificado, alcance y entregables, cronograma con fechas reales
          (Gantt con hitos que coinciden con las entregas), costos, riesgos
          (referencia al capítulo~\ref{cap:riesgos}), calidad, recursos y
          satisfacción de las partes interesadas.
    \item \textbf{Declaración de uso de IA}: qué tareas fueron asistidas y
          cómo se verificaron. Los prompts se guardan en una carpeta
          compartida con la cátedra.
\end{enumerate}
Extensión orientativa: 6 a 10 páginas. Referencias habituales:
\citet{hernandezsampieri2014} para investigación y \citet{pmi2021}
para gestión.
\end{guia}

\section{Introducción}

El presente capítulo describe el marco metodológico adoptado para el
desarrollo de \sistema{}. En línea con los requisitos del proyecto
final, la metodología se articula en dos dimensiones complementarias:
la metodología de investigación, que define cómo se aborda el problema,
se recolecta la información y se validan los resultados; y la
metodología de gestión del proyecto, que establece cómo se organiza,
planifica y ejecuta el trabajo de desarrollo a lo largo del tiempo
disponible. Ambas dimensiones son necesarias y se retroalimentan: el
proceso investigativo informa las decisiones de diseño técnico, mientras
que el marco de gestión garantiza que esas decisiones se implementen de
manera ordenada, dentro del alcance y los plazos establecidos.

\section{Tipo de investigación}

El presente proyecto se enmarca dentro de la \textbf{investigación
aplicada}. A diferencia de la investigación básica --orientada a la
producción de conocimiento teórico sin aplicación inmediata-- la
investigación aplicada parte de un problema real y concreto, y genera
conocimiento con utilidad directa para resolverlo. En el caso de
\sistema{}, el problema es la fragmentación e inaccesibilidad del
mercado de tutorías académicas en Argentina, y el conocimiento generado
se materializa en un artefacto funcional: una plataforma digital que
conecta tutores y alumnos a escala nacional. El aporte original del
proyecto reside en la aplicación de técnicas de \gls{matching} semántico
basadas en \gls{ia} a un dominio educativo local, y en la integración de
funcionalidades de tutoría virtual en un único ecosistema diseñado para
las condiciones geográficas y de conectividad del territorio argentino.
El enfoque es de carácter mixto: predominantemente cualitativo en la
fase de relevamiento del problema y cuantitativo en la fase de
validación.

\section{Metodología de la investigación}

La investigación se estructura siguiendo las etapas del proceso
científico aplicado al desarrollo de sistemas de información.

\subsection{Área temática}

El área temática del proyecto se sitúa en la intersección entre el
desarrollo de plataformas digitales de tipo \gls{marketplace}, la \gls{ia}
aplicada al \gls{matching} de usuarios, y la tecnología educativa
(\gls{edtech}).

\subsection{Planteamiento del problema}

El problema de investigación se formula como la ausencia de una
plataforma nacional que centralice, estandarice y amplíe
geográficamente el acceso al mercado de tutorías particulares en
Argentina.

\subsection{Delimitación de la investigación}

La investigación se delimita a los objetivos específicos OE1 a OE8
definidos en el capítulo~\ref{cap:introduccion}, con foco en el diseño e
implementación del \gls{mvp} de \sistema{} para el mercado argentino.
Quedan fuera del alcance investigativo los mercados regionales, la
integración con organismos estatales y los módulos de accesibilidad y
suscripción premium, los cuales se documentan como trabajo futuro.

\subsection{Marco teórico}

Se realizará una búsqueda sistemática de bibliografía académica y
técnica sobre los temas centrales del proyecto, abarcando plataformas de
\gls{marketplace} educativo, algoritmos de recomendación basados en
\gls{embeddings} semánticos, sistemas de reputación en plataformas
digitales, tecnología \gls{webrtc} aplicada a educación, y el contexto
educativo argentino. Esta revisión se documenta en el
capítulo~\ref{cap:marco-teorico}. Las fuentes a consultar incluyen
artículos científicos, tesis de grado y posgrado, documentación técnica
oficial, informes de mercado de organizaciones como Grand View Research
y Research and Markets, y reportes del Ministerio de Educación de la
Nación Argentina.

\subsection{Diseño concreto de investigación}

A partir de la información relevada, se definirá el diseño técnico del
sistema, abarcando la arquitectura, el modelo de datos, los flujos de
usuario y los contratos de \gls{api}. Esta etapa se documenta en el
capítulo~\ref{cap:marco-tecnologico}.

\subsection{Indicadores}

Los indicadores principales para verificar la hipótesis central se
resumen en la tabla~\ref{tab:indicadores}.

\begin{table}[H]
    \centering
    \caption{Indicadores del proyecto}
    \label{tab:indicadores}
    \small
    \begin{tabularx}{\textwidth}{l L L l}
        \toprule
        \tb{ID} & \tb{Variable} & \tb{Indicador} & \tb{Umbral} \\
        \midrule
        I-01 & Aceptación del \gls{matching} & Tasa de aceptación de las sugerencias del motor por parte del alumno & $> 60\,\%$ de recomendaciones aceptadas en el piloto \\
        I-02 & Calidad percibida de la sesión & Puntuación promedio otorgada por los alumnos en el cuestionario post-sesión (escala 1 a 5) & $\geq 4{,}0$ \\
        I-03 & Cobertura geográfica del piloto & Sesiones completadas entre usuarios de provincias distintas & Al menos 1 sesión exitosa entre dos provincias, incluyendo una del interior \\
        I-04 & Satisfacción y recomendación & \gls{nps} del piloto & Porcentaje positivo (promotores $>$ detractores) \\
        \bottomrule
    \end{tabularx}
\end{table}

\subsection{Técnicas de recolección de datos}

Se implementarán técnicas cualitativas y cuantitativas para la
recolección de datos primarios. La técnica cualitativa consistirá en
entrevistas para validar el problema planteado, identificar puntos de
dolor en el proceso de búsqueda y contratación, y relevar expectativas.
La técnica cuantitativa se basará en encuestas de validación de demanda
para cuantificar la frecuencia de uso de tutorías, los canales actuales
de búsqueda y la disposición a utilizar y pagar por una plataforma
centralizada.

\subsection{Instrumentos de recolección de datos}

Se utilizarán entrevistas semiestructuradas aplicadas a una muestra de
al menos dos tutores particulares activos y cuatro alumnos o familias
que hayan contratado tutorías, complementadas con un formulario breve en
Google Forms orientado a estudiantes universitarios y familias con hijos
en edad escolar, buscando un mínimo de 30 respuestas válidas. El guion
de las entrevistas se transcribe en el apéndice~\ref{apx:entrevista} y
la encuesta completa en el apéndice~\ref{apx:encuesta}.

\subsection{Datos}

Los datos a obtener consistirán en registros de audio provenientes de
las entrevistas semiestructuradas y datos tabulares generados a partir
de las respuestas numéricas y de selección múltiple del formulario
digital.

\subsection{Procesamiento de los datos}

Las entrevistas serán grabadas con el consentimiento del entrevistado
para luego ser transcritas en su totalidad y analizadas mediante
codificación temática, mientras que los datos cuantitativos obtenidos de
los formularios serán agrupados y tabulados para permitir un conteo y
evaluación estadística de las métricas.

\subsection{Análisis de los datos}

Los datos recolectados durante el piloto funcional y la validación de
demanda se analizarán en relación con los indicadores definidos,
presentando los resultados estructurados en el
capítulo~\ref{cap:resultados}.

\subsection{Síntesis y conclusiones}

La síntesis de los datos contrastados con los indicadores de validación
evaluará el cumplimiento de la hipótesis central, presentando las
conclusiones finales y las implicancias del proyecto en el
capítulo~\ref{cap:conclusiones}.

\section{Metodología de gestión del proyecto: Scrum adaptado}

Para la planificación y ejecución del desarrollo de \sistema{} se adopta
\textbf{Scrum adaptado a un contexto unipersonal}, siendo el marco ágil
más apropiado dado que se cuenta con un único desarrollador, un conjunto
de módulos técnicos bien delimitados y una fecha de entrega final fija.
La metodología en cascada fue descartada porque asume que todos los
requisitos pueden definirse con precisión al inicio y que las etapas son
completamente secuenciales, lo cual choca con la necesidad de iteración
en proyectos con \gls{ia} e integraciones complejas como LiveKit o
\gls{api} externas. Kanban fue descartado como marco principal porque,
aunque útil para tareas diarias, no provee la planificación por
objetivos ni los hitos verificables exigidos por el cronograma
académico.

Scrum organiza el trabajo en \textbf{sprints de dos semanas} con
objetivos concretos, revisión del incremento al final de cada ciclo y un
backlog priorizado. En este contexto unipersonal, el desarrollador
asume los roles de Product Owner, Scrum Master y Development Team,
contando con el director del proyecto como stakeholder externo para la
revisión académica. Los artefactos incluyen el Product Backlog, el
Sprint Backlog y el Incremento, mientras que las ceremonias adaptadas
comprenden la Sprint Planning de una hora al inicio de cada ciclo, un
Daily Stand-up diario personal documentado en la bitácora de desarrollo,
la Sprint Review y la Sprint Retrospective.

\section{Gestión del proyecto: dimensiones clave}\label{sec:dimensiones}

En línea con los criterios de evaluación del proyecto y el perfil
directivo en Project Management Professional, se detallan las
dimensiones principales de gestión:

\begin{itemize}
    \item \textbf{Alcance:} definido por los nueve módulos del \gls{mvp}
          descritos en el capítulo~\ref{cap:introduccion}. Requiere
          aprobación explícita cualquier funcionalidad adicional para el
          control del alcance, apoyándose en el uso del Product Backlog,
          el Sprint Backlog y los incrementos funcionales.
    \item \textbf{Tiempos:} el proyecto se estructura en sprints de dos
          semanas a lo largo de seis meses entre junio y noviembre de
          2026, con una dedicación estimada de 15 a 20 horas semanales e
          hitos académicos inamovibles, integrando las ceremonias del
          marco ágil adoptado.
    \item \textbf{Costos:} los costos operativos del \gls{mvp} se estiman
          entre cero y cien dólares mensuales durante el desarrollo,
          cubiertos mediante planes gratuitos y de bajo costo en la nube,
          detallándose plenamente en el capítulo~\ref{cap:economica}.
    \item \textbf{Riesgos:} el análisis detallado se presenta en el
          capítulo~\ref{cap:riesgos}, contemplando factores como demoras
          en la integración de \gls{api} externas, costos de modelos de
          lenguaje superiores a los estimados y la dificultad para
          reclutar usuarios para el piloto.
    \item \textbf{Calidad:} cada módulo desarrollado debe superar pruebas
          unitarias con una cobertura mínima del 70\,\%, además de
          sortear una batería de pruebas de integración previas al
          piloto, midiendo la calidad percibida mediante los indicadores
          específicos de la investigación.
    \item \textbf{Recursos:} un desarrollador full-stack con
          conocimientos en React, Java con Spring Boot, Python y bases de
          datos relacionales, junto al uso de herramientas tecnológicas de
          soporte y la estructura de roles unipersonales descrita.
    \item \textbf{Satisfacción de los interesados:} se medirá a través de
          un sistema de calificaciones post-clase bidireccional para
          mantener altos estándares de calidad, complementándose con el
          análisis del \gls{nps} del grupo piloto para comprobar el grado
          de utilidad percibida y la aplicabilidad de los resultados en
          escenarios reales.
\end{itemize}

\section{Cronograma de entregables}

El cronograma fue construido de atrás hacia adelante tomando como fecha
límite el 28 de noviembre de 2026, fecha estimada de la presentación
final, y distribuyendo los entregables académicos y los sprints de
desarrollo de manera que cada capítulo se apoye en trabajo ya realizado
(tabla~\ref{tab:hitos}).

\begin{table}[H]
    \centering
    \caption{Cronograma de entregables y sprints asociados}
    \label{tab:hitos}
    \small
    \begin{tabularx}{\textwidth}{L L L}
        \toprule
        \tb{Hasta} & \tb{Entregable académico} & \tb{Sprints de desarrollo asociados} \\
        \midrule
        9 jun 2026 & Cap.~2 Metodología & Sprint 0: setup del entorno, estructura del repositorio, wireframes iniciales + spike/prototipo técnico (kill-switch) \\
        28 jul 2026 & Cap.~3 Marco teórico & Sprints 1--2: autenticación y perfiles de usuario (backend + frontend) \\
        18 ago 2026 & Cap.~4 Marco tecnológico & Sprints 3--4: motor de \gls{matching} IA + sistema de búsqueda y filtros \\
        8 sep 2026 & Cap.~5 Justificación económica & Sprints 5--6: reservas + integración MercadoPago \\
        22 sep 2026 & Cap.~6 Análisis de riesgos & Sprint 7: aula virtual (LiveKit + Redis) \\
        13 oct 2026 & Cap.~7 Resultados & Sprints 8--9: módulo IA educativa + sistema de reputación \\
        27 oct 2026 & Cap.~8 Conclusiones & Sprint 10: integración end-to-end, pruebas de sistema y correcciones \\
        28 nov 2026 & Presentación final + producto funcional & Sprint 11: piloto con usuarios reales, análisis de resultados, documentación final \\
        \bottomrule
    \end{tabularx}
\end{table}

Cabe aclarar que Redis en el Sprint 7 cubre el estado propio de la
aplicación --chat en tiempo real, límite de tasa de peticiones-- y no el
estado interno de LiveKit, que al usarse en su modalidad Cloud queda
gestionado del lado del proveedor sin intervención del backend. La
figura~\ref{fig:gantt} muestra la línea de tiempo completa.

\begin{figure}[H]
    \centering
    \begin{ganttchart}[
        hgrid, vgrid={*{6}{draw=none}, dotted},
        x unit=0.42mm,
        y unit chart=6mm, y unit title=7mm,
        time slot format=isodate,
        bar label font=\small, milestone label font=\small,
        title label font=\small,
        bar/.append style={fill=mainColor!60, draw=none},
        milestone/.append style={fill=mainColor, draw=none},
    ]{2026-06-01}{2026-11-28}
        \gantttitlecalendar{month=shortname} \\
        \ganttbar{Sprint 0 --- Setup}{2026-06-01}{2026-06-09} \\
        \ganttbar{Sprints 1--2 --- Autenticación y perfiles}{2026-06-09}{2026-07-28} \\
        \ganttbar{Sprints 3--4 --- Matching IA}{2026-07-28}{2026-08-18} \\
        \ganttbar{Sprints 5--6 --- Reservas y pagos}{2026-08-18}{2026-09-08} \\
        \ganttbar{Sprint 7 --- Aula virtual}{2026-09-08}{2026-09-22} \\
        \ganttbar{Sprints 8--9 --- IA educativa y reputación}{2026-09-22}{2026-10-13} \\
        \ganttbar{Sprint 10 --- Integración}{2026-10-13}{2026-10-27} \\
        \ganttbar{Sprint 11 --- Piloto}{2026-10-27}{2026-11-27} \\
        \ganttmilestone{Entrega final}{2026-11-28}
    \end{ganttchart}
    \caption{Cronograma de sprints del proyecto (ciclo lectivo 2026)}
    \label{fig:gantt}
\end{figure}

\section{Herramientas de gestión y seguimiento}

Para garantizar la trazabilidad del trabajo y la calidad del proceso, se
implementarán diversas herramientas tecnológicas de soporte. Se
utilizará \textbf{GitHub Projects} para la gestión del backlog y los
sprints mediante tableros Kanban integrados al repositorio, donde cada
issue representa una tarea del sprint backlog. El control de versiones
se efectuará en GitHub mediante ramas de características y solicitudes
de revisión antes de la integración a la rama principal. Asimismo, se
mantendrá una bitácora de desarrollo en formato de texto simple
actualizada al cierre de cada jornada para registrar avances y decisiones
técnicas como parte del seguimiento diario.

Para el diseño de interfaz se empleará Figma en la elaboración de
prototipos de alta fidelidad, mientras que Postman servirá para la
documentación y prueba de los endpoints de la \gls{api} REST, y Jest se
configurará como marco de pruebas unitarias del frontend, mientras que
JUnit 5 junto con Mockito cumplen ese rol en el backend Spring Boot,
ambos integrados en el pipeline de integración continua mediante GitHub
Actions. Para el despliegue, el backend se empaqueta como imagen Docker
y se publica mediante Coolify, una plataforma de despliegue autoalojada
que gestiona contenedores, certificados y variables de entorno de forma
equivalente a un PaaS gestionado, alojada primero en infraestructura
propia del desarrollador y migrada a un \gls{vps} antes del piloto.